\chapter{Integración editorial reciente}

\section{Principio de doble registro}
\begin{observation}
Conviene fijar de forma explícita la separación entre tránsito y extensión. Esta separación no reemplaza el doble registro ya introducido en el núcleo; lo precisa en una forma más apta para publicación general.
\end{observation}

\begin{definition}
Sea $X$ un fenómeno legible dentro de una ventana declarada $\Win$. Se introduce la expresión
\[
\Mmeas(X;\Win) = \bigl(n_T(X;\Win), \delta_E(X;\Win)\bigr).
\]
La primera componente cuenta pasos, eventos o transiciones distinguibles. La segunda mide diferencias extensivas dentro del marco declarado.
\end{definition}

\begin{observation}
La cláusula verbal correspondiente debe leerse así: si una diferencia pertenece a la naturaleza del tránsito, se cuenta; si pertenece a la naturaleza de la extensión, se mide; si un fenómeno mezcla ambas naturalezas, su desdoblamiento debe declararse antes de formalizarlo.
\end{observation}

\section{Lectura sobria de D0}
\begin{observation}
D0 no debe cargarse con una metafísica excesiva ni con una magnitud plena prematura. Su sobriedad es una virtud del sistema.
\end{observation}

\begin{definition}
La forma recomendada para su lectura es
\[
\mathfrak{D}_0 = \bigl(\subz, R_{\mathrm{sub}}, (\eps,\Mmeas,\Win), \text{perpetuosidad}\bigr).
\]
\end{definition}

\begin{observation}
Aquí, $\subz$ funciona como símbolo canónico del origen estructural; $R_{\mathrm{sub}}$ como operador de referencia; $(\eps,\Mmeas,\Win)$ como conjunto mínimo de moduladores de legibilidad; y “perpetuosidad” nombra únicamente la condición basal de permanencia del referente.
\end{observation}

\section{Lectura abierta de D1}
\begin{definition}
La formulación editorial mínima de D1 puede fijarse como
\[
(\mathfrak{D}_1,\tau_1,T_1,\psi_1,\text{trayectoria abierta fluctuante, Cronos primitivo}).
\]
\end{definition}

\begin{observation}
En este régimen puede haber deformación, ondulación y torsión local del trayecto, pero no autoencuentro estructural suficiente para producir clausura.
\end{observation}

\begin{observation}
La transición hacia D2 debe declararse por umbral:
\[
\psi_1 \ge \psi^\ast_{\mathrm{cl}} \Longrightarrow \mathfrak{D}_1 \to \mathfrak{D}_2.
\]
Toda transición fuerte debe entrar con criterio, umbral y trazabilidad.
\end{observation}

\section{Línea editorial D0--D7}
\begin{observation}
La secuencia dimensional mínima no debe leerse aquí como tabla escolar de etiquetas, sino como línea de regímenes sucesivos donde el canon permanece fuerte y las lecturas operativas se alojan en él sin sustituirlo.
\end{observation}

\begin{itemize}
\item D0 --- perpetuosidad: referencia basal y ancla de legibilidad.
\item D1 --- Cronos primitivo: trayectoria abierta fluctuante, con dirección y orden pero sin clausura.
\item D2 --- clausura fluida: aparición de autocierre dinámico y propagación límite autocontenida.
\item D3 --- red colmena: coexistencia coherente de trayectorias como tejido relacional.
\item D4 --- enlace o captura resonante: primer régimen no trivial de estabilidad geométrica y coherencia volumétrica.
\item D5 --- persistencia o inercia operativa: continuidad geométrica en banda y estabilidad prolongada.
\item D6 --- memoria: persistencia histórica del trayecto y acumulación estructural.
\item D7 --- confinamiento diferencial: selectores, dominancias y precursores locales de colapso como lecturas operativas admitidas.
\end{itemize}

\section{Cláusula de alcance}
\begin{observation}
MADELEINE no postula que la existencia deba obedecer a nuestras descripciones para realizarse. Postula únicamente que, bajo un marco declarado, ciertas diferencias pueden contarse, ciertas extensiones pueden medirse y ciertas transiciones pueden organizarse de forma trazable hasta el límite que la resolución, la ventana, la unidad escalar relativa y la reproducibilidad efectiva sostengan. Más allá de ese alcance, toda prolongación debe leerse como hipótesis, expansión o reserva.
\end{observation}

